\documentclass[12pt,a4paper]{article}

\usepackage[T1]{fontenc}
\usepackage[utf8]{inputenc}
\usepackage[brazil]{babel}
\usepackage{lmodern}
\usepackage{microtype}
\usepackage{geometry}
\usepackage{setspace}
\usepackage{parskip}
\usepackage{enumitem}
\usepackage{needspace}
\usepackage{amsmath}
\usepackage{amssymb}
\usepackage[hidelinks]{hyperref}

\geometry{left=2.5cm,right=2.5cm,top=2.5cm,bottom=2.5cm}
\setstretch{1.15}
\setlist[enumerate]{leftmargin=*,itemsep=0.45em,topsep=0.35em}
\setlist[itemize]{leftmargin=*,itemsep=0.35em,topsep=0.3em}
\setlength\emergencystretch{4em}

\newcommand{\Section}[1]{\needspace{12\baselineskip}\section{#1}}
\newcommand{\Subsection}[1]{\needspace{10\baselineskip}\subsection{#1}}
\newcommand{\Subsubsection}[1]{\needspace{8\baselineskip}\subsubsection{#1}}

\title{Fundamentos Físicos da Compactação do Solo por Tráfego\\
\large Base conceitual para um protótipo brasileiro tipo Terranimo/Soil2Cover}
\author{}
\date{Data de referência: 03/04/2026}

\begin{document}

\maketitle
\vspace{-1.2em}

%% ====================================================================
\Section{Boussinesq --- propagação de tensão no solo}

%% --------------------------------------------------------------------
\Subsection{O problema original (1885)}

Joseph Boussinesq resolveu analiticamente o seguinte problema: uma \textbf{carga pontual}~$P$ aplicada na superfície de um \textbf{meio semi-infinito, homogêneo, isotrópico e elástico linear}. Qual é a tensão vertical $\sigma_z$ em qualquer ponto $(r,\,z)$ abaixo da superfície?

A solução é:
\begin{equation}
  \sigma_z
  = \frac{3P}{2\pi\, z^2}\;\cos^5\!\alpha
  = \frac{3P}{2\pi\, z^2}\;\frac{1}{\left[1 + \left(\dfrac{r}{z}\right)^{\!2}\right]^{5/2}}
  \label{eq:boussinesq}
\end{equation}
onde $\alpha$ é o ângulo entre a vertical e a linha que liga o ponto de aplicação da carga ao ponto de interesse, $r$ é a distância horizontal e $z$ é a profundidade.

%% --------------------------------------------------------------------
\Subsection{Significado físico}

\begin{itemize}
  \item A tensão decai com o \textbf{quadrado da profundidade} ($\propto 1/z^2$).
  \item Na vertical direta abaixo da carga ($r=0$):
    \begin{equation}
      \sigma_z = \frac{3P}{2\pi\, z^2}
    \end{equation}
  \item Lateralmente, a tensão se espalha formando um \textbf{``bulbo de tensão''} --- as isolinhas formam volumes semelhantes a bulbos de cebola.
  \item A \textbf{profundidade importa mais que a distância lateral}: a 50\,cm de profundidade, a tensão já é uma fração da superficial.
\end{itemize}

%% --------------------------------------------------------------------
\Subsection{Limitações do Boussinesq puro}

\begin{enumerate}
  \item \textbf{Solo não é elástico linear} --- deforma-se plasticamente (irrecuperavelmente) quando a tensão é suficientemente alta.
  \item \textbf{Solo não é homogêneo} --- tem camadas com propriedades distintas (horizonte A, B, C; camada arável vs.\ pé-de-grade).
  \item \textbf{Solo não é isotrópico} --- solos compactados concentram tensão mais verticalmente do que Boussinesq prevê.
  \item \textbf{A carga não é pontual} --- um pneu agrícola tem área de contato de ${\sim}0{,}1$ a $0{,}5$\,m$^2$.
\end{enumerate}

%% ====================================================================
\Section{Correção de Froehlich (1934) --- fator de concentração}

Froehlich introduziu um \textbf{fator de concentração} $\xi$ que generaliza Boussinesq:
\begin{equation}
  \sigma_z = \frac{\xi\, P}{2\pi\, z^2}\;\cos^{\xi+2}\!\alpha
  \label{eq:froehlich}
\end{equation}

\begin{itemize}
  \item $\xi = 3$: recupera Boussinesq exato (solo homogêneo isotrópico).
  \item $\xi > 3$ (ex.: 4, 5, 6): solo mais \textbf{rígido} ou \textbf{pré-compactado} --- a tensão se concentra mais na vertical, penetra mais fundo e se espalha menos lateralmente.
  \item $\xi < 3$ (ex.: 2): solo mais \textbf{mole} --- a tensão se dissipa mais rapidamente e se espalha mais lateralmente.
\end{itemize}

Consequência prática importante: \textbf{um solo já compactado (pé-de-grade, por exemplo) transmite mais tensão para camadas profundas}. O $\xi$ captura esse efeito sem necessidade de MEF.

Na prática, $\xi$ é calibrado por tipo de solo ou estimado a partir da rigidez relativa das camadas. A função \texttt{stressTraffic()} do pacote R \texttt{soilphysics} aceita $\xi$ como parâmetro de entrada.

%% ====================================================================
\Section{TASC --- do pneu real ao campo de tensão (Keller, 2005)}

Boussinesq/Froehlich resolvem para carga pontual. Para um pneu real, o modelo TASC (Tyres/Tracks And Soil Compaction) faz a ponte:

\begin{enumerate}
  \item Calcula a \textbf{área de contato} pneu--solo a partir de: largura do pneu, diâmetro, pressão de inflação e carga na roda.
  \item A forma do contato não é circular --- é uma \textbf{super-elipse} (algo entre elipse e retângulo, controlada por um expoente de forma).
  \item A \textbf{distribuição de pressão} dentro do contato não é uniforme --- é maior no centro e nas bordas dos talões.
  \item O campo de tensão no solo é obtido por \textbf{integração numérica} de Boussinesq/Froehlich sobre essa área de contato.
\end{enumerate}

Em termos práticos, o TASC transforma a especificação de um pneu (ex.: 600/65R28 a 1,2\,bar com 3\,t de carga) em um campo de tensão tridimensional $\sigma(x,\,y,\,z)$ no perfil de solo.

%% ====================================================================
\Section{Pré-adensamento ($\sigma_p$) --- a memória mecânica do solo}

%% --------------------------------------------------------------------
\Subsection{O conceito}

Todo solo possui uma ``memória'' da maior tensão efetiva que já suportou. Essa tensão é o \textbf{pré-adensamento} $\sigma_p$ (ou $p_c$ em algumas notações).

Fisicamente:
\begin{itemize}
  \item Se a tensão aplicada $\sigma < \sigma_p$: o solo se deforma pouco e \textbf{recupera} quando a carga é removida (comportamento elástico, ao longo da linha de recompressão~$\kappa$).
  \item Se $\sigma > \sigma_p$: o solo \textbf{colapsa} --- os grãos se rearranjam, a porosidade diminui irreversivelmente (comportamento plástico, ao longo da linha de compressão virgem~$\lambda$).
\end{itemize}

Analogia: dobrar um clipe de papel. Se se dobra pouco, ele volta à posição original. Se se ultrapassa o limite, ele fica deformado permanentemente. O $\sigma_p$ é esse limite.

%% --------------------------------------------------------------------
\Subsection{Origens do pré-adensamento}

\begin{enumerate}
  \item \textbf{Geológica:} peso de camadas de solo que já existiram acima e foram erodidas.
  \item \textbf{Antrópica (tráfego):} passagens anteriores de máquinas --- a mais relevante neste contexto.
  \item \textbf{Ciclos de umedecimento--secagem:} a sucção matricial gera tensões efetivas altas que ``pré-adensam'' o solo. \textbf{Muito importante nos trópicos.}
  \item \textbf{Cimentação por óxidos:} óxidos de ferro e alumínio colam partículas, gerando resistência estrutural. Específico de solos tropicais (Latossolos).
  \item \textbf{Atividade biológica:} raízes e fauna criam estrutura que resiste ao colapso.
\end{enumerate}

%% --------------------------------------------------------------------
\Subsection{Medição: ensaio oedométrico}

No \textbf{ensaio oedométrico} (ensaio de adensamento), aplica-se carga incremental a uma amostra confinada e mede-se a deformação. No gráfico $e$ (índice de vazios) vs.\ $\log\sigma'$ (tensão efetiva), o ponto de curvatura máxima é o $\sigma_p$ (método de Casagrande).

Limitações práticas:
\begin{itemize}
  \item O ensaio leva dias.
  \item Exige amostra indeformada.
  \item Fornece o valor para \textbf{um ponto, uma profundidade, uma condição de umidade}.
\end{itemize}

%% --------------------------------------------------------------------
\Subsection{PTFs --- Funções de Pedotransferência}

Como não é viável realizar ensaio oedométrico em cada ponto do campo, utilizam-se regressões empíricas:
\begin{equation}
  \sigma_p = f\!\left(\text{argila},\;\text{umidade},\;\rho,\;\text{MO},\;\text{sucção}\right)
  \label{eq:ptf}
\end{equation}

Exemplos já implementados no \texttt{soilphysics}:

\begin{itemize}
  \item \textbf{Horn \& Fleige (2003):} $\sigma_p$ a partir de textura + umidade + densidade. Base: solos europeus temperados.
  \item \textbf{Severiano et al.\ (2013):} $\sigma_p$ para Latossolos brasileiros, função da sucção matricial por faixa de argila. Exemplo para faixa de argila 152\,g/kg:
    \begin{equation}
      \sigma_p = 129\;h^{0{,}15}
    \end{equation}
    onde $h$ é a sucção matricial (kPa).
  \item \textbf{Imhoff et al.\ (2004):} $\sigma_p$ para Hapludox (Latossolo Vermelho):
    \begin{equation}
      \sigma_p = -566{,}764 + 442{,}891\,\rho + 4{,}338\,\text{CC} - 773{,}057\,w
    \end{equation}
    onde $\rho$ é a densidade (Mg/m$^3$), CC é a capacidade de campo (\%) e $w$ é a umidade gravimétrica (g/g).
\end{itemize}

%% ====================================================================
\Section{A lógica do risco --- juntando carga e resistência}

Toda ferramenta tipo Terranimo executa, para cada camada~$z$ do perfil:

\begin{enumerate}
  \item \textbf{Lado da carga:} dado o pneu (TASC) e a integração Boussinesq/Froehlich, calcula-se $\sigma_{\text{aplicada}}(z)$.
  \item \textbf{Lado da resistência:} dada a textura, umidade, densidade e história do solo naquela camada, estima-se $\sigma_p(z)$ via PTF.
  \item \textbf{Decisão:}
    \begin{itemize}
      \item $\sigma_{\text{aplicada}} > \sigma_p$ $\to$ \textbf{vermelho} (compactação adicional ocorrerá).
      \item $\sigma_{\text{aplicada}} \approx \sigma_p$ $\to$ \textbf{amarelo} (risco).
      \item $\sigma_{\text{aplicada}} < \sigma_p$ $\to$ \textbf{verde} (seguro).
    \end{itemize}
\end{enumerate}

O conceito é simples. A dificuldade está nos \textbf{parâmetros de entrada} --- especialmente para solos tropicais.

%% ====================================================================
\Section{Por que o solo tropical é diferente}

%% --------------------------------------------------------------------
\Subsection{Mineralogia}

\begin{itemize}
  \item \textbf{Solos temperados:} dominados por argilominerais 2:1 (montmorillonita, ilita) --- alta CTC, alta plasticidade, comportamento bem previsto pela textura.
  \item \textbf{Solos tropicais:} dominados por argilominerais 1:1 (caulinita) + \textbf{óxidos de ferro} (goethita, hematita) \textbf{e alumínio} (gibbsita).
  \item Os óxidos atuam como \textbf{agentes cimentantes}: colam microagregados formando uma estrutura granular estável mesmo com alto teor de argila.
\end{itemize}

%% --------------------------------------------------------------------
\Subsection{A ``argila que se comporta como areia''}

Um Latossolo Vermelho com 70\% de argila pode apresentar:
\begin{itemize}
  \item \textbf{Comportamento mecânico} similar ao de uma areia: granular, alta permeabilidade, baixa plasticidade.
  \item Isso ocorre porque os óxidos cimentam as partículas de argila em \textbf{microagregados} de 0,05--0,2\,mm que funcionam como grãos de areia.
\end{itemize}

Consequência direta: \textbf{uma PTF europeia que recebe ``70\% argila'' vai prever um solo plástico, coesivo, com alto $\sigma_p$}. Na realidade, esse Latossolo pode ser mais vulnerável à compactação do que o previsto, porque a microestrutura granular colapsa quando a cimentação é rompida.

%% --------------------------------------------------------------------
\Subsection{Ciclos de umedecimento--secagem intensos}

Nos trópicos, a alternância entre estação seca (sucção matricial de centenas de kPa) e estação chuvosa (solo próximo à saturação) gera:

\begin{itemize}
  \item \textbf{Na seca:} sucção alta $\to$ tensão efetiva alta $\to$ pré-adensamento aparente. O solo parece mais resistente.
  \item \textbf{Na chuva:} sucção cai $\to$ perde essa resistência aparente. O solo fica vulnerável.
  \item \textbf{Efeito cumulativo:} os ciclos criam e destroem estrutura, gerando um $\sigma_p$ que \textbf{varia sazonalmente} de forma muito mais dramática que em climas temperados.
\end{itemize}

Isso significa que o $\sigma_p$ não é um número fixo --- depende de \textbf{quando} se está medindo (seca vs.\ chuva). Uma ferramenta embarcada precisa considerar a \textbf{umidade atual} do solo, não um valor tabelado.

%% --------------------------------------------------------------------
\Subsection{Estrutura hierárquica de agregação}

Solos tropicais apresentam uma \textbf{hierarquia de agregação}:

\begin{enumerate}
  \item \textbf{Microagregados} ($< 0{,}25$\,mm) --- cimentados por óxidos, muito estáveis.
  \item \textbf{Macroagregados} ($> 0{,}25$\,mm) --- mantidos por matéria orgânica, raízes e hifas --- mais frágeis.
  \item \textbf{Macroporos} entre agregados --- responsáveis por aeração e drenagem.
\end{enumerate}

O tráfego de máquinas destrói primeiro os \textbf{macroporos} (entre macroagregados). Os microagregados resistem mais. Isso cria um padrão de compactação diferente: perda de macroporosidade sem grande mudança na microporosidade --- algo que um modelo baseado apenas em densidade ($\Delta\rho$) pode não capturar adequadamente.

%% --------------------------------------------------------------------
\Subsection{Implicações para o modelo brasileiro}

Para um ``Terranimo brasileiro'' funcionar, as PTFs precisam considerar:

\begin{enumerate}
  \item \textbf{Não só textura, mas mineralogia} --- teor de Fe$_2$O$_3$, Al$_2$O$_3$, ou pelo menos a classe SiBCS (Latossolo vs.\ Argissolo já carrega informação mineralógica implícita).
  \item \textbf{Sucção matricial atual} como variável de entrada (não apenas umidade gravimétrica).
  \item \textbf{Índice de vazios e grau de saturação} --- capturam melhor o estado do solo tropical que a densidade isolada.
  \item \textbf{Limites de Atterberg} (LL, LP, IP) --- o índice de plasticidade de um Latossolo caulinítico-oxídico é muito diferente de um solo 2:1 com mesma textura.
\end{enumerate}

%% ====================================================================
\Section{Cadeia física para o software proposto}

A cadeia completa, do pneu ao mapa de risco otimizado:

\begin{enumerate}
  \item \textbf{Pneu} (geometria, pressão, carga)
    \begin{itemize}
      \item[$\downarrow$] TASC (Keller, 2005)
    \end{itemize}
  \item \textbf{Área de contato + distribuição de pressão na superfície}
    \begin{itemize}
      \item[$\downarrow$] Boussinesq/Froehlich (integração numérica)
    \end{itemize}
  \item \textbf{Campo de tensão} $\sigma(x,\,y,\,z)$ no perfil
    \begin{itemize}
      \item[$\downarrow$] Comparação camada a camada
    \end{itemize}
  \item \textbf{Resistência do solo:} $\sigma_p(z) = \text{PTF}(\text{SiBCS, textura, óxidos, umidade atual, }\rho)$
    \begin{itemize}
      \item[$\downarrow$]
    \end{itemize}
  \item \textbf{Mapa de risco} (semáforo por camada e por posição GPS)
    \begin{itemize}
      \item[$\downarrow$] Soil2Cover (otimização de rota)
    \end{itemize}
  \item \textbf{Rota otimizada} que minimiza compactação acumulada.
\end{enumerate}

\end{document}
